\chapter{Referencial Teórico}\label{sec:RefTeorico}

\begin{dica}[Dica de metodologia científica: o que escrever no Referencial Teórico]
O referencial teórico apresenta os conceitos e os resultados anteriores de que o
trabalho depende. Não é um resumo de livros: organize por \emph{temas} (e não por
autor), explique só o que o leitor precisa para acompanhar o restante do texto e faça
\emph{síntese crítica} (compare, aponte concordâncias, divergências e lacunas).
\begin{itemize}
  \item Prefira fontes primárias e revisadas por pares: periódicos, anais de eventos
        (\ac{sbc} OpenLib, \ac{acm} Digital Library, \ac{ieee} Xplore), teses e
        dissertações (\ac{bdtd}), livros-texto e normas. Use sites e blogs com critério.
  \item Registre a \emph{busca}: bases consultadas, cadeias de busca, período e
        critérios de inclusão e exclusão. Isso dá transparência e permite repetir a pesquisa.
  \item Parafrasear exige citar a fonte; copiar trecho exige aspas ou recuo e citação.
        Texto de terceiros sem indicação é plágio.
  \item Em \emph{Trabalhos relacionados}, compare os estudos com os \emph{mesmos
        critérios} e termine explicitando o que o seu trabalho acrescenta.
\end{itemize}
\end{dica}

\section{Fundamentação}\label{sec:fundamentacao}

\begin{dica}[Dica de uso da classe: como citar]
\begin{itemize}
  \item \emph{citação indireta} (paráfrase): \texttt{texto \textbackslash cite\{autor2020\}} $\to$ (Autor, 2020);
  \item \emph{no fluxo da frase}: \texttt{\textbackslash textcite\{autor2020\}} $\to$ Autor (2020);
  \item \emph{citação direta curta} (até três linhas): entre aspas, com a página:
        \texttt{\textbackslash enquote\{texto\} \textbackslash cite[p.~12]\{autor2020\}};
  \item \emph{citação direta longa} (mais de três linhas): ambiente \texttt{citacao},
        com recuo de 4 cm, fonte menor e sem aspas;
  \item \emph{citação de citação}: \texttt{\textbackslash textapud\{original\}\{consultado\}};
  \item várias fontes na mesma chamada: \texttt{\textbackslash cite\{a,b,c\}}.
\end{itemize}
O Apêndice~\ref{ap:referencias} mostra o resultado de cada forma de citação.
\end{dica}

Exemplos de citação, em ordem: uma citação indireta, em que se parafraseia a ideia
do autor \cite{autor2020}; uma citação no fluxo da frase, como em \textcite{autor2020};
e uma citação direta curta: segundo o autor, \enquote{o texto transcrito deve vir entre
aspas duplas, seguido da indicação da fonte e da página} \cite[p.~12]{autor2020}.

Quando a citação direta tem mais de três linhas, ela é destacada do texto:

\begin{citacao}
Este é um exemplo de citação direta longa, que tem mais de três linhas. Ela deve ser
destacada com recuo de quatro centímetros em relação à margem esquerda, em letra
menor, em espaço simples e sem aspas. O ponto final encerra a frase e não a citação,
e a indicação da fonte, com a página, vem ao final do trecho \cite[p.~35]{autor2020}.
\end{citacao}

Se for preciso citar uma obra a que não se teve acesso, cite o original por meio da
obra consultada, como em \textapud{memoir}{abntex2classe}; na lista de referências
entra apenas a obra consultada.

\section{Trabalhos Relacionados}\label{sec:TrabRel}

Compare os trabalhos relacionados em um quadro (\autoref{qua:relacionados}). Quadros,
diferentemente de tabelas, têm bordas laterais fechadas e conteúdo predominantemente
textual (Apêndice~\ref{ap:tabelas}).

\begin{quadro}[htbp]
\begin{ilustracao}{\caption{Comparação entre trabalhos relacionados}\label{qua:relacionados}}{\elaboracaopropria}
  \renewcommand{\arraystretch}{1.25}
  \begin{tabularx}{\textwidth}{|>{\raggedright\arraybackslash}p{3.2cm}|X|X|}
    \hline
    \textbf{Trabalho} & \textbf{Contribuição} & \textbf{Limitação} \\
    \hline
    \textcite{autor2020} & Descreva o que o trabalho propõe & Descreva o que ficou em aberto \\
    \hline
    Autor (ano) & Descreva o que o trabalho propõe & Descreva o que ficou em aberto \\
    \hline
    Este trabalho & Descreva a contribuição do seu trabalho & Descreva as limitações do seu trabalho \\
    \hline
  \end{tabularx}
\end{ilustracao}
\end{quadro}
